\documentclass[a4paper,12pt]{article}
\usepackage[T2A]{fontenc}
\usepackage[utf8]{inputenc}
\usepackage[russian]{babel}
\usepackage{amsmath,amssymb,booktabs,tabularx,array,xcolor,tikz,fancyhdr}
\usepackage[margin=23mm,headheight=15pt]{geometry}
\usepackage[colorlinks=true,linkcolor=teal!65!black]{hyperref}
\definecolor{accent}{HTML}{008080}
\setlength{\parindent}{0pt}
\setlength{\parskip}{5pt}
\pagestyle{fancy}
\fancyhf{}
\lhead{\small Объём и площадь поверхности}
\rhead{\small Николь Саркисьянц}
\cfoot{\thepage}
\renewcommand{\headrulewidth}{0.4pt}
\newcommand{\err}{\textcolor{red!70!black}{\textbf{Ошибка:}}}
\newcommand{\idea}{\textcolor{blue!60!black}{\textbf{Ключевая идея:}}}
\newcommand{\result}{\textcolor{teal!60!black}{\textbf{Итог:}}}
\newcommand{\task}[2]{\clearpage\phantomsection\label{task:#1}\section*{Задание № #1. #2}}
\begin{document}
\begin{titlepage}
\centering
\vspace*{30mm}
{\large\color{accent} МАТЕМАТИКА\par}
\vspace{15mm}
{\LARGE Ревью домашней работы\par}
\vspace{9mm}
{\Large\color{accent} Объёмы и площади поверхностей\par многогранников\par}
\vspace{18mm}
{Подготовка к ЕГЭ по профильной математике\par}
\vfill
{Николь Саркисьянц\par}
\vspace{8mm}
{25 ноября 2026 г.\par}
\vspace{16mm}
{\small Лёвин Артём Александрович\par эксклюзивно для Николь Саркисьянц\par}
\vspace{20mm}
\begin{tikzpicture}\draw[accent,line width=.8pt] (0,0) .. controls (2,.8) and (4,-.8) .. (6,0) .. controls (8,.8) and (10,-.8) .. (12,0);\end{tikzpicture}
\end{titlepage}
\section*{Сводная таблица ошибок}
Ответы записаны в порядке $(V;S)$; площади и объёмы выражены в квадратных и кубических единицах соответственно.
\begin{center}
\renewcommand{\arraystretch}{1.5}
\begin{tabularx}{\linewidth}{@{}c p{30mm} X c c@{}}
\toprule № & Тема & Суть ошибки & Ваш ответ & Верный ответ\\\midrule
\hyperref[task:3]{3} & Угловой вырез & Неполный учёт новых граней & $(23;50)$ & $(23;52)$\\
\hyperref[task:5]{5} & Открытый паз & Неверная площадь стенок & $(50;100)$ & $(50;110)$\\
\hyperref[task:10]{10} & Паз и куб & Неполный учёт поверхности паза & $(61;120)$ & $(61;132)$\\\bottomrule
\end{tabularx}
\end{center}
В № 10 ответ в тетради дан в обратном порядке; здесь порядок унифицирован.
\subsection*{Общий принцип}
Для выреза из тела:
\[V=V_0-V_{\text{выреза}},\qquad S=S_0-S_{\text{исчезло}}+S_{\text{появилось}}.\]
При присоединении тела без объёмного перекрытия:
\[V=V_1+V_2,\qquad S=S_1+S_2-2S_{\text{контакта}}.\]
В первом случае учитывайте только участки поверхности, которые действительно исчезают или возникают. Во втором общая грань становится внутренней для обоих тел, поэтому её площадь вычитается дважды.
\clearpage
\thispagestyle{empty}
\vspace*{\fill}\begin{center}{\LARGE\color{accent} Разбор ошибок}\end{center}\vspace*{\fill}
\task{3}{Куб, вырезанный из угла}
\textbf{Условие.} Из угла параллелепипеда $4\times3\times2$ вырезали куб с ребром 1. Удалённый куб имеет три грани на внешней поверхности исходного тела. Найдите объём и площадь полной поверхности оставшегося тела.

\textbf{Ваше решение.} $V_1=24$, $V_2=1$, $V=23$; $S_1=2(12+6+8)=52$. Затем вычтено $2\cdot1=2$, получено $S=50$.

\textbf{Ваш ответ:} $(23;50)$.

\err\ неверная поправка к площади поверхности. Последний правильный шаг: $S_1=52$. Первый неверный шаг --- уменьшение площади на 2 без полного подсчёта исчезнувших и новых граней.

\idea\ в углу сходятся три взаимно перпендикулярные грани. Вырезание куба удаляет три единичных квадрата с внешних граней. Одновременно внутри выреза открываются три новых квадрата той же площади. Все шесть участков принадлежат граням удалённого куба, но выполняют разные роли.

\textbf{Исправление и полное решение.}
\[V_0=4\cdot3\cdot2=24,\quad V_{\text{выреза}}=1^3=1,\quad V=23.\]
\[S_0=2(4\cdot3+3\cdot2+4\cdot2)=52.\]
\[S_{\text{исчезло}}=3\cdot1^2=3,\qquad S_{\text{появилось}}=3\cdot1^2=3.\]
\[S=52-3+3=52.\]
\textbf{Проверка.} Для каждого из трёх направлений одна удалённая квадратная площадка заменяется равной ей параллельной площадкой. Изменение площади равно нулю.

\result\ $V=23$, $S=52$. Объём уменьшился, площадь сохранилась.
\subsection*{Задача на закрепление}
Из угла параллелепипеда $7\times5\times4$ вырезали куб с ребром 2, имеющий три грани на внешней поверхности исходного тела. Найдите объём и площадь поверхности остатка.
\task{5}{Открытый паз между боковыми гранями}
\textbf{Условие.} В параллелепипеде $5\times4\times3$ сделали паз $5\times1\times2$, идущий через всю длину 5, открытый сверху и на двух торцах; по ширине паз расположен строго между боковыми гранями. Найдите объём и площадь полной поверхности.

\textbf{Ваше решение.} $V=60-10=50$, $S_0=94$, $S_-=2(2\cdot1)=4$, $S_+=2(5\cdot1)=10$, $S=94-4+10=100$.

\textbf{Ваш ответ:} $(50;100)$.

\err\ неверно выбраны размеры новых стенок. Последний корректный промежуточный результат --- потеря площади двух торцов, равная 4. Первый неверный шаг --- прибавление 10 в качестве полной поправки от продольных поверхностей. Стенки имеют длину 5 и высоту 2, поэтому их суммарная площадь равна 20. Верхняя полоска также удалена, но её компенсирует дно паза.

\idea\ отдельно учитывайте верхнюю полоску, два торцевых отверстия, дно и две вертикальные стенки паза.

\textbf{Исправление и полное решение.}
\[V=5\cdot4\cdot3-5\cdot1\cdot2=50,\quad S_0=2(20+12+15)=94.\]
\begin{center}
\begin{tabular}{@{}llr@{}}\toprule Участок & Расчёт & Изменение\\\midrule
Верхняя полоска & $5\cdot1$ & $-5$\\
Два торца & $2\cdot1\cdot2$ & $-4$\\
Дно паза & $5\cdot1$ & $+5$\\
Две стенки & $2\cdot5\cdot2$ & $+20$\\\bottomrule\end{tabular}
\end{center}
\[S=94-5-4+5+20=110.\]
\textbf{Проверка.} Поперечное сечение получает две дополнительные вертикальные стороны длины 2; его периметр возрастает на 4. Продольная поверхность увеличивается на $5\cdot4=20$, торцы теряют 4. Чистый прирост равен 16.

\result\ $V=50$, $S=110$.
\subsection*{Задача на закрепление}
В параллелепипеде $8\times6\times5$ вырезали паз $8\times2\times3$, открытый сверху и на обоих торцах, строго между боковыми гранями. Найдите объём и площадь поверхности остатка.
\task{10}{Открытый паз и присоединённый куб}
\textbf{Условие.} Из параллелепипеда $6\times4\times3$ вырезали паз $6\times1\times2$, проходящий через всю длину, открытый сверху и на двух торцах, расположенный строго между боковыми гранями. На свободном участке верхней грани закрепили куб с ребром 1, полностью соприкасающийся с ней нижней гранью и не перекрывающий паз. Найдите объём и площадь полной поверхности фигуры.

\textbf{Ваше решение.} $V=72-12+1=61$, $S_0=108$, $S_-=4$, $S_+=12$, $S_{\text{куба}}=6$. Затем $S=108-4-1+12+6-1=120$.

\textbf{Ваш ответ:} 120; 61 (площадь, объём).

\err\ неверный учёт новых стенок паза. Последний правильный шаг для паза --- потеря площади двух торцов, равная 4. Первый неверный шаг --- $S_+=6\cdot1\cdot2=12$. Это произведение соответствует объёму паза; численно оно также равно площади одной стенки, но стенок две. Прибавление куба с вычитанием двух контактных квадратов выполнено верно.

\idea\ разделите решение на два этапа: сначала поверхность тела с пазом, затем присоединение куба.

\textbf{Исправление и полное решение.}
\[V=6\cdot4\cdot3-6\cdot1\cdot2+1^3=61.\]
\[S_0=2(24+12+18)=108.\]
Верхняя полоска и дно паза имеют одинаковую площадь 6. Два торца теряют $2\cdot1\cdot2=4$. Две стенки добавляют $2\cdot6\cdot2=24$.
\[S_{\text{с пазом}}=108-6-4+6+24=128.\]
При присоединении куба добавляются пять его открытых граней и закрывается единичный квадрат основания:
\[S=128+5\cdot1^2-1^2=132.\]
\textbf{Проверка.} По формуле соединения тел $128+6-2\cdot1=132$. Объём изменяется на $-12+1=-11$, площадь на $-4+24+4=24$.

\result\ $V=61$, $S=132$.
\subsection*{Задача на закрепление}
Из параллелепипеда $8\times6\times4$ вырезали паз $8\times2\times3$, открытый сверху и на обоих торцах, строго между боковыми гранями. На свободном участке верхней грани закрепили куб с ребром 1, прилегающий всей нижней гранью и не перекрывающий паз. Найдите объём и площадь полной поверхности.

\vfill\begin{center}\small Лёвин Артём Александрович\\эксклюзивно для Николь Саркисьянц\end{center}
\end{document}
